\documentclass[aps,prl,onecolumn,superscriptaddress,longbibliography,nofootinbib]{revtex4-2}

\usepackage{amsmath,amssymb,bm}
\usepackage{graphicx}
\usepackage{hyperref}
\usepackage{xcolor}
\usepackage{booktabs}

\newcommand{\br}{\mathbf{r}}
\newcommand{\bv}{\mathbf{v}}
\newcommand{\bx}{\mathbf{X}}
\newcommand{\avg}[1]{\langle #1 \rangle}

\begin{document}

\title{Supplemental Material for:\\
``Superfluid hydrodynamics from correlated many-body dynamics:
Bridging variational Monte Carlo and Gross--Pitaevskii theory
for microwave-shielded Bose gases''}

\author{Xinyuan~Wei}
\affiliation{Department of Physics, University}

\date{\today}

\maketitle

\setcounter{equation}{0}
\setcounter{figure}{0}
\setcounter{table}{0}
\setcounter{section}{0}
\renewcommand{\theequation}{S\arabic{equation}}
\renewcommand{\thefigure}{S\arabic{figure}}
\renewcommand{\thetable}{S\arabic{table}}
\renewcommand{\thesection}{S\arabic{section}}

%======================================================================
\section{Model and interaction potential}
\label{sec:model}
%======================================================================

We consider $N$ identical bosons of mass $m$ in a cylindrically symmetric
harmonic trap with Hamiltonian
\begin{equation}
\hat{H}(t) = \sum_{i=1}^N\!\left[
-\frac{\hbar^2}{2m}\nabla_i^2 + V_{\mathrm{tr}}(\br_i,t)\right]
+ \sum_{i<j} v(\br_{ij}),
\label{eq:H}
\end{equation}
where the trap potential is
\begin{equation}
V_{\mathrm{tr}}(\br,t) = \frac{1}{2}m\left[\omega_\rho(t)^2\rho^2
+ \omega_z^2 z^2\right],
\end{equation}
with $\rho^2 = x^2 + y^2$ the cylindrical radial coordinate.  We work
in natural units $\hbar=m=1$ and measure lengths in units of the
oscillator length $a_{\mathrm{ho}}=\sqrt{\hbar/(m\omega_0)}$ with
$\omega_0=1$.

The two-body interaction models microwave-shielded dipolar
molecules~\cite{Karman2018}:
\begin{equation}
v(\br) = V_{\mathrm{dip}}(r,\theta) + V_{\mathrm{shield}}(r,\theta),
\end{equation}
where $r=|\br|$ and $\theta$ is the polar angle with respect to the
$z$-axis.  The dipolar term is
\begin{equation}
V_{\mathrm{dip}} = \frac{C_3}{r^3}(3\cos^2\theta - 1),
\end{equation}
and the repulsive shielding term is
\begin{equation}
V_{\mathrm{shield}} = \frac{C_6}{r^6}(\sin^4\theta
+ w_{12}\sin^2\theta\cos^2\theta),
\end{equation}
with $C_3=10^{-3}$, $C_6=10^{-6}$, and $w_{12}=2.0$ in oscillator
units.  Both potentials are softened at short range by
$r\to\sqrt{r^2+\epsilon_r^2}$ with $\epsilon_r=10^{-3}$ to avoid
numerical singularities.

The system is initialized in the ground state of an isotropic trap
$\omega_\rho=\omega_z=1.0$.  At $t=0$, the radial trap frequency is
instantaneously quenched to $\omega_\rho=2.0$ while $\omega_z$ remains
unchanged, exciting a radial breathing mode.

%======================================================================
\section{Jastrow--Feenberg wavefunction: full specification}
\label{sec:wavefunction}
%======================================================================

\subsection{One-body field $u_1$}

The one-body field is parametrized with a separable radial-axial
structure:
\begin{equation}
u_1(\rho,z,t) = u_\rho(\rho,t) + u_z(z,t),
\end{equation}
where each component consists of a harmonic reference plus Gaussian
radial basis function corrections:
\begin{align}
u_\rho(\rho,t) &= -\frac{1}{2}(\omega_{\mathrm{ref}}+a_\rho(t))\rho^2
+ \sum_{k=1}^{n_\rho}c_{\rho,k}(t)\,g_0(\rho;\ell_{\rho,k}),
\label{eq:u1_rho}\\
u_z(z,t) &= -\frac{1}{2}(\omega_{\mathrm{ref}}+a_z(t))z^2
+ \sum_{k=1}^{n_z}c_{z,k}(t)\,g_0(z;\ell_{z,k}),
\label{eq:u1_z}
\end{align}
with $\omega_{\mathrm{ref}}=1.0$ and the Gaussian basis functions
\begin{equation}
g_0(r;\ell) = \exp\!\left(-\frac{r^2}{2\ell^2}\right).
\end{equation}

The complex variational parameters $\{a_\rho(t), a_z(t),
c_{\rho,1}(t),\ldots,c_{\rho,n_\rho}(t), c_{z,1}(t),\ldots,
c_{z,n_z}(t)\}$ are time-dependent.  Their real parts control the
density profile and their imaginary parts encode the velocity field.

The length scales are distributed logarithmically:
\begin{equation}
\ell_k = \ell_{\min}\left(\frac{\ell_{\max}}{\ell_{\min}}\right)^{(k-1)/(n-1)},
\quad k=1,\ldots,n,
\end{equation}
with $\ell_{\min}=0.1$, $\ell_{\max}=3.0$, and $n_\rho=n_z=6$.
This gives a total of $n_{u_1}=2+n_\rho+n_z=14$ complex one-body
parameters.

\subsection{Two-body Jastrow $u_2$}

The two-body factor is expanded in partial waves up to $\ell=2$:
\begin{equation}
u_2(r,\cos\theta) = f_0(r) + f_2(r)\,P_2(\cos\theta),
\end{equation}
where $P_2(x)=(3x^2-1)/2$ is the Legendre polynomial.  The radial
components are
\begin{align}
f_0(r) &= u_{\mathrm{cusp}}(r)
+ \sum_{k=1}^{n_2} c_{0,k}\,g_0(r;\ell_{2,k}),
\label{eq:f0}\\
f_2(r) &= \sum_{k=1}^{n_2} c_{2,k}\,g_2(r;\ell_{2,k}),
\label{eq:f2}
\end{align}
with $g_2(r;\ell)=r^2\exp(-r^2/(2\ell^2))$ a regular $\ell=2$ basis
function and $n_2=10$ basis functions with length scales
$\ell_2\in[0.01,0.7]$.

\subsection{Cusp treatment}

The short-range cusp function captures the two-body scattering boundary
condition imposed by the $C_6/r^6$ repulsion:
\begin{equation}
u_{\mathrm{cusp}}(r) = -\frac{\kappa}{\sqrt{r^2+\epsilon_c^2}}\,
\exp\!\left(-\frac{r^2}{2r_c^2}\right),
\label{eq:cusp}
\end{equation}
with cusp strength $\kappa=0.1$, softening parameter
$\epsilon_c=5\times10^{-3}$, and cutoff radius $r_c=0.01$.
The $1/r$ singularity is regularized by $\epsilon_c$ to avoid
numerical overflow, while the Gaussian envelope ensures that the cusp
affects only the short-range pair structure ($r\lesssim r_c$).  These
parameters are fixed (not variational) and determined by the scattering
properties of the two-body potential.

The total variational space has $n_{u_1}+2n_2=14+20=34$ parameters,
of which $n_{u_1}=14$ one-body parameters are time-dependent and
$2n_2=20$ two-body parameters are fixed in the frozen-$u_2$
approximation.

%======================================================================
\section{Time-dependent variational principle: detailed derivation}
\label{sec:tdvp}
%======================================================================

\subsection{General TDVP formulation}

The TDVP projects the Schr\"odinger equation onto the variational
tangent space spanned by derivatives of $\ln\Psi$ with respect to the
variational parameters $\{\theta_\alpha\}$.  Starting from the
Schr\"odinger equation $i\partial_t\Psi = \hat{H}\Psi$ and the
variational ansatz $\Psi(\bx;\bm\theta(t))$, the McLachlan variational
principle~\cite{McLachlan1964} minimizes
\begin{equation}
\delta\!\left\|\,i\partial_t\Psi - \hat{H}\Psi\,\right\|^2 = 0
\end{equation}
with respect to $\dot{\bm\theta}$.  This yields the equations of
motion~\cite{Carleo2012}:
\begin{equation}
\sum_\beta S_{\alpha\beta}\,\dot\theta_\beta = -i\,F_\alpha,
\label{eq:tdvp}
\end{equation}
where the overlap matrix (quantum geometric tensor) and force vector are
\begin{align}
S_{\alpha\beta} &= \avg{\Delta O_\alpha^*\,\Delta O_\beta}_t
= \avg{O_\alpha^* O_\beta}_t - \avg{O_\alpha^*}_t\avg{O_\beta}_t,
\label{eq:S_def}\\
F_\alpha &= \avg{\Delta O_\alpha^*\,\Delta E_L}_t
= \avg{O_\alpha^* E_L}_t - \avg{O_\alpha^*}_t\avg{E_L}_t,
\label{eq:F_def}
\end{align}
with logarithmic derivatives $O_\alpha =
\partial_{\theta_\alpha}\ln\Psi$.

\subsection{Local energy}

For the Jastrow ansatz $\Psi=e^U$ with
$U=\sum_i u_1(\br_i,t)+\sum_{i<j}u_2(r_{ij})$, the local energy is
\begin{equation}
E_L = \frac{\hat{H}\Psi}{\Psi} = \sum_{i=1}^N\left[
-\frac{1}{2}(\nabla_i^2 U + |\nabla_i U|^2) + V_{\mathrm{tr}}(\br_i,t)
\right] + \sum_{i<j} v(\br_{ij}),
\label{eq:EL}
\end{equation}
where
\begin{align}
\nabla_i U &= \nabla u_1(\br_i,t) + \sum_{j\neq i}\nabla_i u_2(r_{ij}),\\
\nabla_i^2 U &= \nabla^2 u_1(\br_i,t)
+ \sum_{j\neq i}\nabla_i^2 u_2(r_{ij}).
\end{align}

\subsection{Monte Carlo estimation}

The expectation values in Eqs.~\eqref{eq:S_def}--\eqref{eq:F_def} are
estimated by sampling configurations $\bx^{(s)}$ from $|\Psi(\bx,t)|^2$
using the Metropolis algorithm.  We employ single-particle moves:
for each walker and each particle $i$, a displacement
$\Delta\br\sim\mathrm{Uniform}(-\delta/2,\delta/2)^3$ is proposed
with step size $\delta=0.45$, and accepted with probability
$\min(1,|\Psi_{\mathrm{new}}|^2/|\Psi_{\mathrm{old}}|^2)$.

We use $N_w=65{,}536$ parallel walkers (GPU-accelerated using
JAX~\cite{Bradbury2018}), with $n_{\mathrm{burn}}=30$ equilibration
sweeps followed by $n_{\mathrm{meas}}=10$ measurement sweeps separated
by 3 thinning sweeps, per RK4 stage.

The TDVP linear system~\eqref{eq:tdvp} is solved via eigendecomposition
of $S$:
\begin{equation}
\dot\theta_\alpha = -i\sum_\beta S^{-1}_{\alpha\beta}F_\beta,
\end{equation}
with eigenvalue truncation at $\lambda_{\min}/\lambda_{\max}=10^{-4}$
and a diagonal shift $S\to S+\eta\mathbb{1}$ with $\eta=10^{-2}$.

For imaginary-time propagation (ground-state preparation), we use
$\dot\theta = -F/S$ with $d\tau=10^{-3}$ for 2000 steps.  For
real-time evolution after the quench, the step size is
$dt=2\times10^{-4}$ with adaptive time-stepping that limits the RMS
parameter change to $\|\dot\theta\,dt\|_{\mathrm{rms}}<0.01$.

%======================================================================
\section{Frozen-$u_2$ approximation: detailed derivation of
hydrodynamic equations}
\label{sec:frozen_derivation}
%======================================================================

\subsection{Motivation: stiffness of $u_2$}

Our tVMC simulations reveal that during the trap quench
$\omega_\rho:1\to 2$, the two-body Jastrow parameters $(c_{0,k},
c_{2,k})$ change by less than 1\% of their ground-state values, while
the one-body parameters $(a_\rho, a_z, c_{\rho,k}, c_{z,k})$ undergo
order-unity changes.  This reflects the physical expectation that
short-range pair correlations equilibrate much faster than the global
density redistribution driven by the trap quench.

\subsection{TDVP with continuous one-body field}

Freezing $u_2$ at its ground-state value (real, time-independent) and
promoting $u_1(\br,t)$ to a continuous variational field, the
logarithmic derivative becomes the microscopic density operator:
\begin{equation}
O(\br;\bx) = \frac{\delta\ln\Psi}{\delta u_1(\br,t)}
= \sum_{i=1}^N\delta^{(3)}(\br-\br_i) \equiv \hat\rho(\br).
\end{equation}

The TDVP equation~\eqref{eq:tdvp} becomes a continuum integral equation:
\begin{equation}
\int d^3r'\;S(\br,\br';t)\,\partial_t u_1(\br',t) = -i\,F(\br;t),
\label{eq:tdvp_cont}
\end{equation}
where
\begin{align}
S(\br,\br';t) &= \avg{\Delta\hat\rho(\br)\,\Delta\hat\rho(\br')}_t,\\
F(\br;t) &= \avg{\Delta\hat\rho(\br)\,\Delta E_L}_t.
\end{align}

The density--density correlator can be decomposed as
\begin{equation}
S(\br,\br') = \rho(\br)\delta^{(3)}(\br-\br')
+ \rho^{(2)}(\br,\br') - \rho(\br)\rho(\br'),
\end{equation}
where $\rho(\br)=\avg{\hat\rho(\br)}_t$ is the one-body density and
$\rho^{(2)}(\br,\br')$ is the two-body density.

\subsection{Decomposition into amplitude and phase}

We decompose the complex one-body field:
\begin{equation}
u_1(\br,t) = a(\br,t) + i\theta(\br,t), \qquad a,\theta\in\mathbb{R}.
\end{equation}
Because $u_2$ is real and frozen, the many-body wavefunction has the
structure
\begin{equation}
\Psi = e^{\sum_i a(\br_i) + i\sum_i\theta(\br_i)
+ \sum_{i<j}u_2(r_{ij})},
\end{equation}
so the many-body phase is purely one-body: $\arg\Psi = \sum_i
\theta(\br_i)$.  The flow is therefore irrotational with velocity
\begin{equation}
\bv(\br,t) = \nabla\theta(\br,t).
\end{equation}
No backflow terms arise because $u_2$ carries no time-dependent phase.

\subsection{Internal energy functional}

Define the internal energy as the expectation value of kinetic plus
interaction energy in the frozen-$u_2$ state at density $\rho$:
\begin{equation}
E_{\mathrm{int}}[\rho;u_2] = \avg{\hat{T}+\hat{V}_{\mathrm{int}}}_t.
\end{equation}
If $E_{\mathrm{int}}$ depends on $\rho$ only through the local density
(LDA), then $E_{\mathrm{int}}^{\mathrm{LDA}} = \int\rho(\br)
\varepsilon(\rho(\br))\,d^3r$, where $\varepsilon(\rho)$ is the energy
per particle of the uniform system.  The TDVP action then takes the form
\begin{equation}
\mathcal{S} = \int\!dt\,d^3r\left[
\rho\,\partial_t\theta - \frac{1}{2}\rho|\nabla\theta|^2
- \rho\,V_{\mathrm{tr}} - \rho\,\varepsilon(\rho)\right],
\end{equation}
and the stationarity conditions $\delta\mathcal{S}/\delta\theta=0$
and $\delta\mathcal{S}/\delta a=0$ (where $\rho$ depends on $a$) yield
the closed hydrodynamic system:
\begin{align}
\text{Continuity:}\quad
&\partial_t\rho + \nabla\!\cdot(\rho\bv) = 0,
\label{eq:continuity}\\
\text{Euler:}\quad
&\partial_t\bv + \nabla\!\left(\frac{1}{2}v^2 + V_{\mathrm{tr}}
+ \mu(\rho)\right) = 0,
\label{eq:euler}
\end{align}
where the chemical potential is
\begin{equation}
\mu(\rho) = \frac{d[\rho\varepsilon(\rho)]}{d\rho}
= \varepsilon(\rho) + \rho\varepsilon'(\rho).
\end{equation}

\subsection{Gross--Pitaevskii reformulation}

Equations~\eqref{eq:continuity}--\eqref{eq:euler} can be rewritten as
a single GPE for $\psi=\sqrt\rho\,e^{i\theta}$:
\begin{equation}
i\partial_t\psi = \left[-\frac{1}{2}\nabla^2
+ V_{\mathrm{tr}} + \mu(|\psi|^2)\right]\psi,
\label{eq:gpe}
\end{equation}
where the Laplacian provides both the quantum pressure
($Q=-\frac{1}{2}\nabla^2\sqrt\rho/\sqrt\rho$) and the flow kinetic
energy.  This formulation avoids the $\rho\to 0$ singularity of the
Madelung variables and provides a single complex equation.

\subsection{Quantum pressure}

The quantum pressure term arises from the density dependence of the
overlap kernel $S(\br,\br')$ beyond the LDA.  In the GPE
formulation~\eqref{eq:gpe}, it is automatically included through the
$-\frac{1}{2}\nabla^2\psi$ term.  For the weakly interacting $N=8$
system ($E_{\mathrm{int}}/E_{\mathrm{tot}}\approx 3.5\%$), quantum
pressure provides the dominant repulsive mechanism.  Without it, the
dynamics is qualitatively incorrect.  For $N=128$
($E_{\mathrm{int}}/E_{\mathrm{tot}}\approx 21\%$), both quantum
pressure and the EOS contribute significantly.

%======================================================================
\section{Equation of state: VMC procedure and fitting}
\label{sec:eos}
%======================================================================

\subsection{Uniform-system VMC}

To extract the EOS, we perform VMC calculations for a uniform system
of $N$ particles in a cubic box with periodic boundary conditions at
a grid of densities.  The two-body Jastrow $u_2(r,\hat\br)$ is frozen
at the ground-state values from the trapped-system tVMC, while $u_1$
is taken as constant.

At each density, we use $N_w=8{,}192$ walkers with 200 equilibration
sweeps and 500 measurement sweeps (thinning by 5) to estimate the
energy per particle $\varepsilon(\rho)$.  Error bars are obtained from
a blocking analysis.

For $N=8$, we scan 20 densities logarithmically over
$\rho\in[0.05,3.0]$.  For $N=128$, the higher peak densities
encountered during compression require extending the range to
$\rho\in[0.05,15.0]$ with 30 density points.

\subsection{Polynomial fitting}

We fit $\varepsilon(\rho)$ to a polynomial using weighted least-squares:
\begin{equation}
\varepsilon(\rho) = a_1\rho + a_2\rho^2 + a_3\rho^3.
\end{equation}
Model selection between quadratic ($a_3=0$) and cubic is performed
using the Akaike information criterion (AIC).

For $N=8$, the cubic fit is preferred:
\begin{equation}
a_1 = 0.102,\quad a_2 = 0.023,\quad a_3 = -0.007,
\end{equation}
with reduced $\chi^2_\nu=0.41$.  The negative $a_3$ introduces a
stability limit $d\mu/d\rho=0$ at $\rho_{\mathrm{stab}}\approx 2.67$,
above which we regularize by linear extrapolation.

For $N=128$, a quadratic suffices:
\begin{equation}
a_1 = 0.111,\quad a_2 = 0.001,\quad a_3 = 0,
\end{equation}
which is unconditionally stable ($d\mu/d\rho>0$ for all $\rho$).

\subsection{Emergent effective interactions}

The chemical potential $\mu(\rho)=2a_1\rho+3a_2\rho^2+\cdots$ admits
a physically transparent interpretation.  Comparing with the standard
GPE nonlinearity $\mu=g\rho$:
\begin{itemize}
\item The leading term $2a_1\rho$ defines a \emph{renormalized
two-body coupling} $g_{\mathrm{eff}}=2a_1$ that absorbs the effect of
frozen pair correlations into an effective contact interaction.
\item The next-order term $3a_2\rho^2$ has the structure of an
\emph{effective three-body interaction} $g_3|\psi|^4$ with
$g_3=3a_2$.
\end{itemize}
This three-body term is not a genuine microscopic three-body force but
emerges from integrating out the density-dependent pair correlations
encoded in $u_2$.  The resulting GPE
\begin{equation}
i\partial_t\psi = \left[-\frac{1}{2}\nabla^2 + V_{\mathrm{tr}}
+ g_{\mathrm{eff}}|\psi|^2 + g_3|\psi|^4\right]\psi
\end{equation}
provides a concrete, non-perturbative realization of the extended GPE
discussed in the context of quantum droplets~\cite{Petrov2015} and
dipolar condensates~\cite{Chomaz2023}.

%======================================================================
\section{Hydrodynamic PDE solver: numerical details}
\label{sec:solver}
%======================================================================

\subsection{Cylindrical coordinates}

Exploiting azimuthal symmetry, the 3D GPE reduces to 2D in cylindrical
coordinates $(r,z)$:
\begin{equation}
\nabla^2\psi = \frac{\partial^2\psi}{\partial r^2}
+ \frac{1}{r}\frac{\partial\psi}{\partial r}
+ \frac{\partial^2\psi}{\partial z^2}.
\end{equation}

For $N=8$: domain $r\in[0,5]$, $z\in[-5,5]$, grid $128\times 256$.

For $N=128$: domain $r\in[0,6]$, $z\in[-8,8]$, grid $192\times 512$.

\subsection{Boundary conditions}

At $r=0$: azimuthal symmetry requires $\partial_r\psi|_{r=0}=0$,
implemented via a ghost point $\psi(-\Delta r)=\psi(\Delta r)$.  The
$(1/r)\partial_r\psi$ term is evaluated using L'H\^opital's rule:
$\lim_{r\to 0}(1/r)\partial_r\psi = \partial_r^2\psi|_{r=0}$.

At outer boundaries ($r=R_{\max}$, $z=\pm Z_{\max}$): Dirichlet
$\psi=0$.

\subsection{Fourth-order finite differences}

Interior derivatives use fourth-order central differences:
\begin{align}
f''_i &= \frac{-f_{i-2}+16f_{i-1}-30f_i+16f_{i+1}-f_{i+2}}
{12\Delta x^2},\\
f'_i &= \frac{-f_{i+2}+8f_{i+1}-8f_{i-1}+f_{i-2}}{12\Delta x}.
\end{align}
Near boundaries, ghost-point values are determined by symmetry or
Dirichlet conditions.  The fourth-order scheme is critical for
frequency accuracy: the second-order scheme gives
$\omega_{\mathrm{br}}=3.99$, while fourth-order gives $4.19$
(matching tVMC).

\subsection{Time integration}

RK4 with dispersive CFL condition:
\begin{equation}
\Delta t \leq C_{\mathrm{CFL}}\frac{(\min\Delta x)^2}{\pi},
\quad C_{\mathrm{CFL}}=0.3.
\end{equation}
An additional constraint $\Delta t\leq 2C_{\mathrm{CFL}}/|\mu_{\max}|$
limits phase accumulation per step.

\subsection{Imaginary-time ground state}

The GPE ground state is obtained by imaginary-time propagation:
\begin{equation}
\psi \to \psi - d\tau\,\hat{H}_{\mathrm{GPE}}\psi,\qquad
\psi \to \psi\sqrt{N/\|\psi\|^2},
\end{equation}
with $d\tau = 0.3\Delta x^2/\pi$, iterated until $\delta\mu<10^{-8}$.
Starting from a Gaussian guess, convergence is achieved in
$\sim 30{,}000$ iterations for $N=8$ and $\sim 45{,}000$ for $N=128$.

\subsection{Observables}

At regular intervals we compute (using cylindrical trapezoidal
integration $\int f\,2\pi r\,dr\,dz$):
\begin{align}
N(t) &= \int|\psi|^2\,2\pi r\,dr\,dz,\\
\avg{\rho^2}(t) &= \frac{1}{N}\int r^2|\psi|^2\,2\pi r\,dr\,dz,\\
\avg{z^2}(t) &= \frac{1}{N}\int z^2|\psi|^2\,2\pi r\,dr\,dz,\\
E(t) &= \int\!\left[-\frac{1}{2}\psi^*\nabla^2\psi
+ V|\psi|^2 + |\psi|^2\varepsilon(|\psi|^2)\right]
2\pi r\,dr\,dz.
\end{align}
The kinetic energy uses the same discrete Laplacian as the dynamics,
ensuring consistent energy conservation.

%======================================================================
\section{Quantitative comparison: tVMC vs GPE}
\label{sec:tvmc_vs_gpe}
%======================================================================

Table~\ref{tab:results} summarizes the comparison between frozen-$u_2$
tVMC and the GPE solver for both $N=8$ and $N=128$ after the quench
$\omega_\rho:1\to 2$.  For tVMC, we report the energy conservation as
the drift in the stochastic energy estimator after the quench; the
instantaneous fluctuation is $\sim 5\%$ for $N=8$ and $\sim 7\%$ for
$N=128$ due to Monte Carlo noise, but the systematic drift is
negligible.

\begin{table}[h]
\caption{Comparison of frozen-$u_2$ tVMC and GPE (imaginary-time ground
state, fourth-order Laplacian) for the quench $\omega_\rho:1\to 2$.
$\delta E/E$ for tVMC is the energy drift (not stochastic fluctuation);
for GPE it is the deterministic conservation error.}
\label{tab:results}
\begin{ruledtabular}
\begin{tabular}{lcccc}
 & \multicolumn{2}{c}{$N=8$} & \multicolumn{2}{c}{$N=128$} \\
\cline{2-3}\cline{4-5}
 & tVMC & GPE & tVMC & GPE \\
\hline
$E_{\mathrm{gs}}$
& 12.42 & 12.44 & 276.8 & 269.0 \\
$\langle\rho^2\rangle_0$
& 1.051 & 1.058 & 1.555 & 1.561 \\
$\langle z^2\rangle_0$
& 0.525 & 0.529 & 0.772 & 0.780 \\
$\omega_{\mathrm{br}}$
& 4.19 & 4.19 & 4.19 & 4.19 \\
$\langle\rho^2\rangle$ range
& [0.26,\,1.06] & [0.26,\,1.06] & [0.37,\,1.56] & [0.37,\,1.56] \\
$\langle z^2\rangle$ range
& [0.52,\,0.61] & [0.53,\,0.61] & [0.59,\,2.05] & [0.56,\,2.00] \\
$|\delta E/E|$
& $<0.01\%$ & $0.1\%$ & $2.8\%$ & $0.05\%$ \\
$\max|\delta N/N|$
& --- & $7\!\times\!10^{-4}$ & --- & $3\!\times\!10^{-4}$ \\
\end{tabular}
\end{ruledtabular}
\end{table}

%======================================================================
\section{Convergence tests}
\label{sec:convergence}
%======================================================================

\subsection{Spatial convergence}

Table~\ref{tab:convergence} shows conservation diagnostics for $N=8$
at different grid resolutions and Laplacian orders.

\begin{table}[h]
\caption{Conservation diagnostics for the GPE simulation ($N=8$,
$t\in[0,3]$).}
\label{tab:convergence}
\begin{ruledtabular}
\begin{tabular}{lccc}
Grid & $\nabla^2$ order & $\max|\delta N/N|$ & $\max|\delta E/E|$\\
\hline
$128\times 256$ & 2nd & $1.09\times10^{-3}$ & $1.20\times10^{-3}$\\
$128\times 256$ & 4th & $7.30\times10^{-4}$ & $1.29\times10^{-3}$\\
$160\times 320$ & 4th & $4.65\times10^{-4}$ & $8.22\times10^{-4}$\\
\end{tabular}
\end{ruledtabular}
\end{table}

\subsection{Temporal convergence}

Halving the CFL factor from 0.3 to 0.15 (doubling the time steps)
produces identical conservation to machine precision, confirming that
the RK4 temporal error is negligible and the conservation errors are
purely spatial.

\subsection{Conservation behavior}

The conservation errors are \emph{oscillatory}: they peak during
maximum compression (steepest density gradients) and nearly vanish
when the cloud re-expands.  At $t=1.5$ (full expansion for $N=8$),
$|\delta N/N|=3.2\times10^{-6}$ and $|\delta E/E|=2.8\times10^{-6}$.

For $N=128$, the larger grid ($192\times 512$) achieves
$\max|\delta N/N|=3.1\times10^{-4}$ and
$\max|\delta E/E|=4.5\times10^{-4}$ over $t\in[0,6]$.

%======================================================================
\section{Scaling ansatz}
\label{sec:scaling}
%======================================================================

As a further benchmark, we compare with the Castin--Dum scaling
ansatz~\cite{CastinDum1996,Kagan1996}.  Assuming self-similar evolution
$\rho(\br,t) = (\lambda_r^2\lambda_z)^{-1}
\rho_0(r/\lambda_r,z/\lambda_z)$, the scaling parameters satisfy
\begin{align}
\ddot\lambda_r &= -\omega_\rho^2\lambda_r
+ \frac{\omega_{\rho,0}^2}{\lambda_r^3\lambda_z},\\
\ddot\lambda_z &= -\omega_z^2\lambda_z
+ \frac{\omega_{z,0}^2}{\lambda_r^2\lambda_z^2}.
\end{align}
This is the non-interacting ($\gamma=1$) limit, appropriate for
benchmarking the trap-dominated regime.  The scaling ansatz predicts
a breathing frequency close to $2\omega_\rho$ but overestimates the
oscillation amplitude because it neglects quantum pressure corrections
to the self-similar assumption.

%======================================================================
\section{Comparison between frozen-$u_2$ tVMC and full tVMC}
\label{sec:full_vs_frozen}
%======================================================================

For $N=8$, we verify that the frozen-$u_2$ approximation is accurate
by comparing with a full tVMC simulation where all 34 parameters
evolve.  Table~\ref{tab:full_frozen} summarizes the comparison.

\begin{table}[h]
\caption{Comparison of full and frozen-$u_2$ tVMC for $N=8$.}
\label{tab:full_frozen}
\begin{ruledtabular}
\begin{tabular}{lcc}
& Full tVMC & Frozen-$u_2$ tVMC\\
\hline
$E_{\mathrm{gs}}$ & 12.433 & 12.423\\
$\avg{\rho^2}_0$ & 1.049 & 1.051\\
$\avg{z^2}_0$ & 0.525 & 0.525\\
$\avg{\rho^2}$ range & [0.264,\,1.060] & [0.264,\,1.055]\\
$\avg{z^2}$ range & [0.522,\,0.613] & [0.522,\,0.613]\\
$\omega_{\mathrm{br}}$ (dominant) & 3.14 & 4.19\\
\end{tabular}
\end{ruledtabular}
\end{table}

The ground state and oscillation amplitudes agree to within 0.5\%.
The breathing frequency differs: the full tVMC shows a dominant mode
at $\omega\approx 3.14$ due to additional mode coupling through the
time-dependent $u_2$, while the frozen-$u_2$ tVMC and GPE both give
$\omega\approx 4.19\approx 2\omega_\rho$, consistent with the
hydrodynamic prediction for this weakly interacting regime.

For $N=128$, only the full tVMC simulation was performed (the
frozen-$u_2$ tVMC was not run separately).  The GPE with the
VMC-derived EOS reproduces the full tVMC breathing frequency
($\omega=4.19$) and oscillation amplitudes to within 2\%, as
discussed in the main text.

%======================================================================
\section{Partial-wave convergence: $\ell=4$ analysis}
\label{sec:l4}
%======================================================================

The microwave-shielded interaction [Eq.~(2) of the main text] contains
two angular structures: the dipolar term $\propto 3\cos^2\!\theta-1 =
2P_2(\cos\theta)$ and the shielding term $\propto \sin^4\!\theta$.
While the dipolar term couples only $\ell=0$ and $\ell=2$, the
$\sin^4\!\theta$ term contains a $P_4$ component and in principle
requires $\ell=4$ in the Jastrow expansion:
\begin{equation}
u_2(r,\cos\theta) = f_0(r) + f_2(r)\,P_2(\cos\theta)
+ f_4(r)\,P_4(\cos\theta).
\end{equation}

To test whether $\ell=4$ matters, we performed a ground-state
optimization for $N=8$ with $n_4=6$ additional $\ell=4$ basis functions
$g_4(r;\ell)=r^4\exp(-r^2/(2\ell^2))$.
Table~\ref{tab:l4} compares the ground-state properties.

\begin{table}[h]
\caption{Ground-state comparison with and without $\ell=4$ ($N=8$).}
\label{tab:l4}
\begin{ruledtabular}
\begin{tabular}{lcc}
 & $\ell\leq 2$ (v7) & $\ell\leq 4$ (v8) \\
\hline
$E_{\mathrm{gs}}$ & 12.423 & 12.439 \\
$\avg{\rho^2}_0$ & 1.051 & 1.051 \\
$\avg{z^2}_0$ & 0.525 & 0.525 \\
$|c_0|_{\mathrm{rms}}$ & 1.503 & 1.503 \\
$|c_2|_{\mathrm{rms}}$ & 0.046 & 0.046 \\
$|c_4|_{\mathrm{rms}}$ & --- & 0.001 \\
$|c_4/c_0|$ & --- & 0.07\% \\
\end{tabular}
\end{ruledtabular}
\end{table}

The $\ell=4$ contribution is negligible: $|c_4|_{\mathrm{rms}}$ is
0.07\% of $|c_0|_{\mathrm{rms}}$, and even the $\ell=2$ anisotropy is
only 3\% of $\ell=0$.  The ground-state energy and moments are
unaffected to within the stochastic error of the VMC.

\subsection{Dynamical stability of $\ell=4$ parameters}

During real-time quench dynamics, the weakly constrained $c_4$
parameters are susceptible to numerical instability in the TDVP.
With standard regularization ($\eta=10^{-4}$), $|c_4/c_0|$ grows
from 0.07\% to $\sim 39\%$ over $t\in[0,3]$, despite negligible
physical coupling.  Two diagnostic tests confirm this is a numerical
artifact:

\begin{enumerate}
\item \emph{No-quench test}: With $\omega_\rho$ held at 1.0 (no
perturbation), $|c_4/c_0|$ still grows from 0.02\% to 0.48\%---a
monotonic random walk that should be identically zero for a stationary
state.
\item \emph{Regularization sensitivity}: Increasing the diagonal shift
from $\eta=10^{-4}$ to $\eta=10^{-3}$ suppresses the growth to 1.4\%,
confirming noise-driven amplification of weakly constrained directions.
\end{enumerate}

These tests justify our choice to retain only $\ell\leq 2$ in the
production simulations.  For the frozen-$u_2$ approximation used to
derive the hydrodynamic equations, this instability is entirely avoided
since the pair parameters are not evolved.

%======================================================================
\section{EOS fitting: extended range for $N=128$}
\label{sec:eos_extended}
%======================================================================

The $N=128$ system reaches peak densities $\rho_{\mathrm{peak}}\approx
29$ during maximum radial compression.  The EOS must therefore be
reliable over a much wider density range than for $N=8$.

We extended the uniform VMC scan to $\rho_{\max}=15$ with 30 density
points.  The results are well described by a quadratic:
$\varepsilon(\rho)=0.111\rho+0.001\rho^2$.  The EOS is nearly linear,
reflecting the dominance of the long-range dipolar interaction at high
density.

We compared both quadratic (poly2) and cubic (poly3) fits.  The poly3
fit has a stability limit at $\rho_{\mathrm{stab}}\approx 10.7$, beyond
which $d\mu/d\rho<0$.  Since the simulation reaches $\rho\approx 29$,
this requires regularization.  The poly2 fit is unconditionally stable
and gives slightly better agreement with tVMC---particularly in the
axial dynamics, where $\avg{z^2}_{\max}=1.94$ (poly2) vs.\ $1.84$
(poly3) compared to tVMC $2.05$.

The small quadratic coefficient ($a_2/a_1\approx 0.007$) means the
effective three-body coupling $g_3=3a_2\approx 0.002$ is weak relative
to the two-body coupling $g_{\mathrm{eff}}=2a_1\approx 0.22$.  The
beyond-mean-field correction is therefore small but non-negligible,
contributing $\sim 1\%$ of the chemical potential at typical
densities.

\bibliography{letter_refs}

\end{document}
